\documentclass[11pt, a4paper]{article}

% bfw-style.tex — gemeinsame Style-Definitionen für alle Handouts/Aufgabenhefte
% Voraussetzung: Kompilation mit xelatex (wegen fontspec)

\usepackage[ngerman]{babel}
% Babel-ngerman macht " zu einem aktiven Shorthand (z.B. für "a -> ä). In unseren
% Texten verwenden wir " als normales Zeichen — daher Shorthand komplett ausschalten.
\shorthandoff{"}
\usepackage{geometry}
\geometry{a4paper, top=2.4cm, bottom=2.4cm, left=2.3cm, right=2.3cm,
          headheight=15pt, headsep=0.7cm, footskip=1.1cm}

\usepackage{fontspec}
% Fallback-Kette: Source Sans Pro -> Calibri -> Segoe UI -> Latin Modern Sans
% (Latin Modern ist immer mit MiKTeX vorhanden)
\IfFontExistsTF{Source Sans Pro}{%
  \setmainfont{Source Sans Pro}[Ligatures=TeX]%
}{\IfFontExistsTF{Calibri}{%
  \setmainfont{Calibri}[Ligatures=TeX]%
}{\IfFontExistsTF{Segoe UI}{%
  \setmainfont{Segoe UI}[Ligatures=TeX]%
}{%
  \setmainfont{Latin Modern Sans}[Ligatures=TeX]%
}}}
\IfFontExistsTF{Source Code Pro}{%
  \setmonofont{Source Code Pro}[Scale=0.92]%
}{\IfFontExistsTF{Consolas}{%
  \setmonofont{Consolas}[Scale=0.92]%
}{%
  \setmonofont{Latin Modern Mono}[Scale=0.92]%
}}

% Gesperrte Versalien für Labels/Titelbalken (Kapitälchen-Ersatz, funktioniert
% mit jeder OpenType-Schrift — Latin Modern Sans hat keine echten Kapitälchen)
\newcommand{\bfwlabelfont}{\footnotesize\bfseries\addfontfeatures{LetterSpace=8.0}}

\usepackage{xcolor}
% Farbpalette: hell, freundlich, modern
\definecolor{bfwPrimary}{HTML}{0EA5A4}    % Petrol/Türkis - Primär
\definecolor{bfwPrimaryDark}{HTML}{0F766E} % dunkler Petrol für Headings
\definecolor{bfwAccent}{HTML}{F59E0B}     % Amber - Akzent
\definecolor{bfwSuccess}{HTML}{10B981}    % Grün - Erfolg / Lösung
\definecolor{bfwWarn}{HTML}{F97316}       % Orange - Achtung
\definecolor{bfwInk}{HTML}{0F172A}        % fast Schwarz
\definecolor{bfwInkSoft}{HTML}{475569}    % gedämpftes Grau
\definecolor{bfwBgSoft}{HTML}{F8FAFC}     % off-white
\definecolor{bfwBgTeal}{HTML}{ECFEFF}     % helles Türkis
\definecolor{bfwBgAmber}{HTML}{FEF3C7}    % helles Gelb
\definecolor{bfwBgRose}{HTML}{FFE4E6}     % helles Rosé für Eselsbrücken
\definecolor{bfwTabHead}{HTML}{E4F3F2}    % zarte Kopfzeilen-Hinterlegung für Tabellen

\usepackage[colorlinks=true, linkcolor=bfwPrimaryDark, urlcolor=bfwPrimaryDark]{hyperref}
\usepackage{lastpage}
\usepackage{enumitem}
\setlist[itemize]{leftmargin=*, itemsep=2pt, topsep=2pt}
\setlist[enumerate]{leftmargin=*, itemsep=2pt, topsep=2pt}

\usepackage[most]{tcolorbox}
\usepackage{booktabs}
\usepackage{colortbl}
\usepackage{tikz}
\usetikzlibrary{decorations.pathmorphing, positioning, shapes.geometric, arrows.meta}

% ---------------------------------------------------------------------------
% Einheitliches Tabellen-Design
%   - etwas mehr Zeilenluft, dezente graue Linien (wirkt auch mit \hline ruhiger)
%   - Kopfzeile einer Tabelle bei Bedarf zart hinterlegen: \rowcolor{bfwTabHead}
% ---------------------------------------------------------------------------
\renewcommand{\arraystretch}{1.25}
\arrayrulecolor{bfwInkSoft!50}
\setlength{\heavyrulewidth}{1pt}
\setlength{\lightrulewidth}{0.5pt}

% ---------------------------------------------------------------------------
% Boxen — klare farbige Titelbalken mit gesperrten Versal-Labels
% (Titel bleibt per Option überschreibbar: \begin{merksatz}[title={...}])
% ---------------------------------------------------------------------------
\tcbset{bfwbox/.style={%
  enhanced, breakable,
  boxrule=0.5pt, arc=2.5pt,
  left=10pt, right=10pt, top=8pt, bottom=8pt,
  toptitle=4pt, bottomtitle=4pt,
  titlerule=0pt,
  fonttitle=\bfwlabelfont,
  coltitle=white
}}

% Merkkästchen — wichtige Definition / Kernpunkt
\newtcolorbox{merksatz}[1][]{%
  bfwbox,
  colback=bfwBgTeal,
  colframe=bfwPrimaryDark,
  colbacktitle=bfwPrimaryDark,
  title={MERKE},
  #1
}

% Eselsbrücke — Mnemoniks
\newtcolorbox{eselsbruecke}[1][]{%
  bfwbox,
  colback=bfwBgRose,
  colframe=bfwAccent!85!black,
  colbacktitle=bfwAccent!85!black,
  title={ESELSBRÜCKE},
  #1
}

% Beispiel-Box
\newtcolorbox{beispiel}[1][]{%
  bfwbox,
  colback=bfwBgSoft,
  colframe=bfwInkSoft,
  colbacktitle=bfwInkSoft,
  title={BEISPIEL},
  #1
}

% Lösungs-Box (für Aufgabenhefte)
\newtcolorbox{loesung}[1][]{%
  bfwbox,
  colback=bfwSuccess!7,
  colframe=bfwSuccess!65!black,
  colbacktitle=bfwSuccess!65!black,
  title={LÖSUNG},
  #1
}

% Achtung-Box — typische Fehler / Stolperfallen
\newtcolorbox{achtung}[1][]{%
  bfwbox,
  colback=bfwWarn!6,
  colframe=bfwWarn!85!black,
  colbacktitle=bfwWarn!85!black,
  title={ACHTUNG},
  #1
}

% Formel-Box — für zentrale Formeln (groß, ruhig, ohne Titelbalken)
\newtcolorbox{formelbox}[1][]{%
  enhanced,
  colback=bfwBgTeal,
  colframe=bfwPrimaryDark,
  boxrule=1pt, arc=3pt,
  left=10pt, right=10pt, top=6pt, bottom=6pt,
  #1
}

% Glossar-Eintrag
\newcommand{\glossareintrag}[2]{%
  \par\noindent\textbf{\color{bfwPrimaryDark}#1} \enspace #2\par\smallskip
}

% ---------------------------------------------------------------------------
% Abschnitts-Design: farbiges Nummern-Quadrat + feine Linie unter \section,
% dezent farbige \subsection/\subsubsection
% ---------------------------------------------------------------------------
\usepackage{titlesec}
\newcommand{\bfwSecBadge}[1]{%
  \begin{tikzpicture}[baseline={([yshift=-0.62em]n.center)}]
    \node[fill=bfwPrimaryDark, text=white, font=\normalsize\bfseries,
          minimum width=1.55em, minimum height=1.55em, inner sep=1pt,
          rounded corners=1.5pt] (n) {#1};
  \end{tikzpicture}}
\titleformat{\section}
  {\Large\bfseries\color{bfwPrimaryDark}}
  {\bfwSecBadge{\thesection}}{0.6em}{}
  [\vspace{-0.2em}{\color{bfwPrimary!60}\titlerule[0.8pt]}]
\titleformat{\subsection}
  {\large\bfseries\color{bfwPrimaryDark}}
  {\textcolor{bfwPrimary}{\thesubsection}}{0.5em}{}
\titleformat{\subsubsection}
  {\normalsize\bfseries\color{bfwInk}}
  {\textcolor{bfwPrimary}{\thesubsubsection}}{0.4em}{}
\titlespacing*{\section}{0pt}{1.6em}{1em}
\titlespacing*{\subsection}{0pt}{1.2em}{0.5em}

% TikZ-Stil "skizzenhaft" — etwas weicher als Rough.js, aber stabil
\tikzset{
  roughbox/.style = {
    draw=bfwPrimaryDark, thick, fill=bfwBgTeal,
    rounded corners=4pt, inner sep=6pt
  },
  roughaccent/.style = {
    draw=bfwAccent, thick, fill=bfwBgAmber,
    rounded corners=4pt, inner sep=6pt
  },
  roughline/.style = {
    draw=bfwInkSoft, thick, -{Stealth[length=2.5mm]}
  }
}

% Bullet-Symbole (bewusst kein FontAwesome — vermeidet Font-Installations-Probleme)
\newcommand{\faLightbulb}{$\bullet$}
\newcommand{\faBookmark}{$\bullet$}

% ---------------------------------------------------------------------------
% Kopf-/Fußzeile
%   Kopf links:  Wortlogo „Wirtschaftsfachwirt"
%   Kopf rechts: Dokument-Kurztext, gesetzt via \kopfzeile{...}
%   Fuß links:   © Can Siebert · 2026 — Fuß rechts: Seite X von Y
%   Titelseite (Seite 1) bleibt ohne Kopf-/Fußzeile (\handouttitel setzt
%   \thispagestyle{empty}).
% ---------------------------------------------------------------------------
\usepackage{fancyhdr}
\newcommand{\bfwKopfzeileText}{}
\newcommand{\kopfzeile}[1]{\renewcommand{\bfwKopfzeileText}{#1}}
\pagestyle{fancy}
\fancyhf{}
\fancyhead[L]{\small\bfseries\color{bfwPrimaryDark}Wirtschaftsfachwirt}
\fancyhead[R]{\small\color{bfwInkSoft}\bfwKopfzeileText}
\renewcommand{\headrulewidth}{0.6pt}
\renewcommand{\headrule}{{\color{bfwPrimary}\hrule width\headwidth height 0.6pt}}
\fancyfoot[L]{\small\color{bfwInkSoft}\copyright\ Can Siebert\,\textperiodcentered\,2026}
\fancyfoot[R]{\small\color{bfwInkSoft}Seite \thepage\ von \pageref*{LastPage}}
% Auch \thispagestyle{plain} (z.B. durch \maketitle) soll leer bleiben:
\fancypagestyle{plain}{\fancyhf{}\renewcommand{\headrulewidth}{0pt}\renewcommand{\headrule}{}}

% ---------------------------------------------------------------------------
% \handouttitel{Obertitel}{Untertitel}{Kurs-Label}{Termin-Zeile}
%   Einheitlicher professioneller Titelblock: Dachzeile, farbige Headline,
%   Untertitel, farbige Trennlinie und dezente Meta-Zeile (Kurs/Termin/Dozent).
%   Ersetzt \title/\maketitle. Direkt nach \begin{document} aufrufen.
% ---------------------------------------------------------------------------
\newcommand{\handouttitel}[4]{%
  \thispagestyle{empty}%
  \begingroup
  \setlength{\parindent}{0pt}%
  {\bfwlabelfont\color{bfwPrimary}WIRTSCHAFTSFACHWIRT (IHK)\par}
  \vspace{0.55em}
  {\huge\bfseries\color{bfwPrimaryDark}#1\par}
  \vspace{0.5em}
  {\large\color{bfwInkSoft}#2\par}
  \vspace{0.9em}
  {\color{bfwPrimary}\rule{\linewidth}{2pt}}\par
  \vspace{0.55em}
  {\small\color{bfwInkSoft}#3\ \,\textperiodcentered\,\ #4\ \,\textperiodcentered\,\ Dozent: Can Siebert\par}
  \endgroup
  \vspace{1.4em}%
}


\title{\vspace*{-1.5cm}\Huge\bfseries\color{bfwPrimaryDark}Aufgabenheft Lektion~1\\[0.4em]
  \large\color{bfwInkSoft}\textnormal{Volkswirtschaftliche Grundbegriffe \& Preisbildung --- Übungen im IHK-Klausurformat}}
\author{\textsf{Wirtschaftsfachwirt --- 1.1 Volkswirtschaftliche Grundlagen \textperiodcentered{} \small\copyright\ Can Siebert\,\textperiodcentered\,2026}}
\date{Selbstlernmaterial}

\begin{document}
\maketitle
\thispagestyle{empty}

\begin{merksatz}[title={\bfseries So nutzen Sie dieses Aufgabenheft}]
Bearbeiten Sie alle Aufgaben zunächst \emph{ohne} in die Lösungen zu schauen. Schreiben
Sie Ihre Antworten in vollständigen Sätzen --- die IHK-Prüfung verlangt formulierte
Begründungen, nicht bloße Stichworte. Beachten Sie die \emph{Operatoren} (nennen,
erläutern, beurteilen, begründen) --- sie geben den geforderten Antwort-Tiefegrad vor.
Erst danach öffnen Sie den Lösungsteil ab Seite~\pageref{loesungen} und vergleichen.
Ergänzend: Übungsband Volkswirtschaft, Kapitel~1 (Übungen 1--17).
\end{merksatz}

\tableofcontents
\vspace{1em}
\hrule

\section{Aufgaben}

\subsection*{Aufgabe~1 --- Ökonomisches Prinzip (8 Punkte)}

Die Marketingleiterin eines Möbelherstellers sagt: „Wir wollen mit unserem festen
Werbebudget von 50.000~Euro so viele Kunden wie möglich erreichen.`` Ihr Kollege aus der
Produktion erwidert: „Und ich muss die geplanten 1.000 Schränke mit möglichst geringen
Kosten fertigen.``

\begin{enumerate}[label=(\alph*)]
  \item \textbf{Erklären} Sie, warum Menschen und Unternehmen überhaupt wirtschaften müssen. \hfill (3~P.)
  \item \textbf{Ordnen} Sie die beiden Aussagen jeweils dem Maximal- bzw. Minimalprinzip zu und \textbf{begründen} Sie Ihre Zuordnung. \hfill (4~P.)
  \item \textbf{Beurteilen} Sie die Aussage: „Wirtschaftlich handelt, wer mit minimalen Mitteln maximalen Erfolg erzielt.`` \hfill (1~P.)
\end{enumerate}

\subsection*{Aufgabe~2 --- Grundbegriffe (10 Punkte)}

\begin{enumerate}[label=(\alph*)]
  \item \textbf{Unterscheiden} Sie die Begriffe Bedürfnis, Bedarf und Nachfrage anhand eines selbst gewählten Beispiels. \hfill (4~P.)
  \item \textbf{Definieren} Sie den Begriff Wirtschaftssubjekt und \textbf{nennen} Sie drei Beispiele. \hfill (3~P.)
  \item \textbf{Nennen} Sie die drei klassischen volkswirtschaftlichen Produktionsfaktoren sowie den zunehmend als eigenständig betrachteten vierten Faktor. \hfill (3~P.)
\end{enumerate}

\subsection*{Aufgabe~3 --- Marktarten und vollkommener Markt (12 Punkte)}

\begin{enumerate}[label=(\alph*)]
  \item \textbf{Definieren} Sie den Begriff „Markt``. \hfill (2~P.)
  \item \textbf{Nennen} Sie die drei Marktarten nach dem Marktgegenstand mit je einem Beispiel. \hfill (3~P.)
  \item \textbf{Nennen} Sie die fünf Bedingungen des vollkommenen Marktes. \hfill (5~P.)
  \item \textbf{Begründen} Sie, warum die Börse dem vollkommenen Markt sehr nahe kommt. \hfill (2~P.)
\end{enumerate}

\subsection*{Aufgabe~4 --- Marktformen (10 Punkte)}

Ordnen Sie die folgenden Situationen jeweils der zutreffenden Marktform zu und
\textbf{begründen} Sie kurz:

\begin{enumerate}[label=(\alph*)]
  \item Wenige Mineralölkonzerne beliefern Millionen Autofahrer mit Kraftstoff. \hfill (2~P.)
  \item Ein regionaler Wasserversorger beliefert alle Haushalte einer Stadt. \hfill (2~P.)
  \item Viele Händler verkaufen auf Wochenmärkten an viele Kunden. \hfill (2~P.)
  \item Der Staat fragt als einziger Straßenbauleistungen bei vielen Bauunternehmen nach. \hfill (2~P.)
  \item Die Bundeswehr kauft Kampfflugzeuge bei einem einzigen Hersteller in Deutschland. \hfill (2~P.)
\end{enumerate}

\subsection*{Aufgabe~5 --- Bestimmungsfaktoren und Kurvenverschiebungen (14 Punkte)}

Der Absatz von Wärmepumpen in Deutschland steigt deutlich, nachdem die Preise für Heizöl
und Erdgas stark gestiegen sind. Gleichzeitig sinken durch technischen Fortschritt die
Produktionskosten der Hersteller.

\begin{enumerate}[label=(\alph*)]
  \item \textbf{Erläutern} Sie, um welche Güterbeziehung es sich zwischen Erdgasheizungen und Wärmepumpen handelt. \hfill (3~P.)
  \item \textbf{Erläutern} Sie, wie sich der Anstieg der Gaspreise auf die Nachfragekurve für Wärmepumpen auswirkt (Richtung der Verschiebung mit Begründung). \hfill (4~P.)
  \item \textbf{Erläutern} Sie, wie sich die gesunkenen Produktionskosten auf die Angebotskurve auswirken. \hfill (3~P.)
  \item \textbf{Unterscheiden} Sie allgemein die „Bewegung auf der Nachfragekurve`` von der „Verschiebung der Nachfragekurve``. \hfill (4~P.)
\end{enumerate}

\subsection*{Aufgabe~6 --- Preiselastizität (10 Punkte)}

Der Preis für ein Medikament steigt um 20~\%, die nachgefragte Menge sinkt daraufhin um
2~\%. Bei einem Modeartikel führt eine Preiserhöhung um 10~\% zu einem Nachfragerückgang
um 25~\%.

\begin{enumerate}[label=(\alph*)]
  \item \textbf{Berechnen} Sie die direkte Preiselastizität der Nachfrage für beide Güter. \hfill (4~P.)
  \item \textbf{Interpretieren} Sie beide Ergebnisse (elastisch/unelastisch) und \textbf{begründen} Sie den Unterschied. \hfill (4~P.)
  \item \textbf{Erklären} Sie den Begriff Kreuzpreiselastizität. \hfill (2~P.)
\end{enumerate}

\subsection*{Aufgabe~7 --- Preisbildung im Gleichgewicht (16 Punkte)}

Für einen Markt bei vollständiger Konkurrenz gilt folgende Marktübersicht:

\begin{center}
\begin{tabular}{c c c}
\toprule
\textbf{Preis (€)} & \textbf{Angebotene Menge (Stück)} & \textbf{Nachgefragte Menge (Stück)}\\
\midrule
10 & 1.000 & 200\\
8  & 800   & 400\\
6  & 600   & 600\\
4  & 400   & 800\\
2  & 200   & 1.000\\
\bottomrule
\end{tabular}
\end{center}

\begin{enumerate}[label=(\alph*)]
  \item \textbf{Bestimmen} Sie Gleichgewichtspreis und Gleichgewichtsmenge. \hfill (2~P.)
  \item \textbf{Beschreiben} Sie die Marktsituation bei einem Preis von 10~€ und die daraus folgende Anpassungsreaktion (Begriff + Mechanismus). \hfill (5~P.)
  \item \textbf{Beschreiben} Sie die Marktsituation bei einem Preis von 2~€ (Begriff: Käufer- oder Verkäufermarkt?) und die Anpassungsreaktion. \hfill (5~P.)
  \item \textbf{Erklären} Sie die Begriffe Konsumentenrente und Produzentenrente. \hfill (4~P.)
\end{enumerate}

\subsection*{Aufgabe~8 --- Funktionen des Marktpreises (8 Punkte)}

\textbf{Nennen und erläutern} Sie die vier Funktionen des Marktpreises mit je einem
Beispiel. \hfill (8~P.)

\subsection*{Aufgabe~9 --- Preisbildung bei unvollständiger Konkurrenz (12 Punkte)}

\begin{enumerate}[label=(\alph*)]
  \item \textbf{Erklären} Sie, wie ein Anbieter im unvollkommenen Polypol einen „monopolistischen Absatzbereich`` erlangt und welche Preisspielräume sich daraus ergeben. \hfill (5~P.)
  \item \textbf{Erläutern} Sie die Begriffe Preisführerschaft und Parallelverhalten im Oligopol am Beispiel des Tankstellenmarktes. \hfill (4~P.)
  \item \textbf{Begründen} Sie, warum ein Angebotsmonopolist als „Preisfixierer`` bezeichnet wird, und \textbf{nennen} Sie eine Grenze seiner Preissetzungsmacht. \hfill (3~P.)
\end{enumerate}

\newpage
\section{Lösungen}\label{loesungen}

\begin{loesung}[title={Lösung Aufgabe~1 --- Ökonomisches Prinzip}]
\textbf{(a)} Die menschlichen Bedürfnisse sind nahezu unbegrenzt, die zu ihrer Befriedigung
verfügbaren Güter und Mittel jedoch knapp. Aus diesem Spannungsverhältnis (Knappheit) folgt
die Notwendigkeit, mit den vorhandenen Mitteln planvoll umzugehen --- also zu wirtschaften.

\textbf{(b)} Marketingleiterin: \emph{Maximalprinzip} --- die Mittel (50.000~€) sind fest
vorgegeben, der Erfolg (Kundenreichweite) soll maximiert werden. Produktionskollege:
\emph{Minimalprinzip} --- das Ziel (1.000 Schränke) ist fest vorgegeben, der Mitteleinsatz
(Kosten) soll minimiert werden.

\textbf{(c)} Die Aussage ist ökonomisch nicht haltbar: Minimale Mittel und maximaler Erfolg
können nicht gleichzeitig fixiert bzw. optimiert werden. Entweder sind die Mittel gegeben
(Maximalprinzip) oder das Ziel (Minimalprinzip).
\end{loesung}

\begin{loesung}[title={Lösung Aufgabe~2 --- Grundbegriffe}]
\textbf{(a)} \emph{Bedürfnis:} empfundener Mangel (z.\,B. Wunsch nach einem neuen Auto).
\emph{Bedarf:} das mit Kaufkraft ausgestattete Bedürfnis (ich habe 30.000~€ zur Verfügung).
\emph{Nachfrage:} der am Markt wirksam werdende Bedarf (ich gehe zum Händler und will
kaufen).

\textbf{(b)} Ein Wirtschaftssubjekt ist das kleinste wirtschaftende Grundelement, das
selbstständig wirtschaftliche Entscheidungen treffen kann. Beispiele: Privathaushalt,
Unternehmen, staatliche Körperschaft.

\textbf{(c)} Arbeit, Boden, Kapital; zunehmend außerdem Wissen (Know-how) als vierter
Produktionsfaktor.
\end{loesung}

\begin{loesung}[title={Lösung Aufgabe~3 --- Marktarten und vollkommener Markt}]
\textbf{(a)} Der Markt ist das Zusammentreffen von Angebot und Nachfrage, bei dem sich der
Preis bildet.

\textbf{(b)} Faktormarkt (z.\,B. Arbeitsmarkt), Gütermarkt (z.\,B. Konsumgüter),
Geld- und Kapitalmarkt/Finanzmarkt (z.\,B. Aktienmarkt).

\textbf{(c)} (1) Homogenität der Güter, (2) Fehlen persönlicher, räumlicher und zeitlicher
Präferenzen, (3) vollkommene Markttransparenz, (4) unendlich schnelle Reaktion der
Marktteilnehmer (Punktmarkt), (5) polypolistische Marktstruktur mit freiem Marktzutritt.

\textbf{(d)} An der Börse sind die gehandelten Güter (z.\,B. Aktien einer Gesellschaft)
homogen, die Bedingungen des Punktmarktes sind erfüllt und die Makler besitzen
Markttransparenz; lediglich Beschränkungen beim Marktzutritt bestehen.
\end{loesung}

\begin{loesung}[title={Lösung Aufgabe~4 --- Marktformen}]
\textbf{(a)} Angebotsoligopol --- wenige Anbieter, viele Nachfrager.
\textbf{(b)} Angebotsmonopol --- ein Anbieter, viele Nachfrager.
\textbf{(c)} (Zweiseitiges) Polypol --- viele Anbieter, viele Nachfrager.
\textbf{(d)} Nachfragemonopol --- ein Nachfrager (Staat), viele Anbieter.
\textbf{(e)} (Zweiseitiges) Monopol --- ein Nachfrager (Bundeswehr), ein Anbieter.
\end{loesung}

\begin{loesung}[title={Lösung Aufgabe~5 --- Bestimmungsfaktoren und Kurvenverschiebungen}]
\textbf{(a)} Es handelt sich um \emph{Substitutionsgüter} (sich ersetzende Güter): Die
Wärmepumpe kann die Gasheizung als Heizsystem ersetzen.

\textbf{(b)} Steigt der Preis des Substituts (Gas/Heizöl), weichen Nachfrager auf
Wärmepumpen aus. Die Nachfrage nach Wärmepumpen steigt \emph{bei jedem Preis} ---
die Nachfragekurve verschiebt sich nach \emph{rechts} (Nachfrageerhöhung), obwohl sich
der Preis der Wärmepumpen selbst nicht geändert hat.

\textbf{(c)} Sinkende Produktionskosten verbessern die Gewinnaussichten; die Anbieter
weiten das Angebot aus. Die Angebotskurve verschiebt sich nach \emph{rechts}
(Angebotserhöhung); tendenziell sinkt der Gleichgewichtspreis, die Menge steigt.

\textbf{(d)} \emph{Bewegung auf der Kurve:} Der Preis des Gutes selbst ändert sich, die
nachgefragte Menge passt sich entlang der Kurve an. \emph{Verschiebung der Kurve:} Andere
Einflussgrößen als der Preis ändern sich (Einkommen, Preise anderer Güter,
Bedürfnisstruktur, Erwartungen) --- die gesamte Kurve wandert nach links oder rechts.
\end{loesung}

\begin{loesung}[title={Lösung Aufgabe~6 --- Preiselastizität}]
\textbf{(a)} Medikament: $\varepsilon = \frac{-2\,\%}{+20\,\%} = -0{,}1$ (Betrag: 0,1).
Modeartikel: $\varepsilon = \frac{-25\,\%}{+10\,\%} = -2{,}5$ (Betrag: 2,5).

\textbf{(b)} Medikament: $|\varepsilon| < 1$ $\Rightarrow$ \emph{unelastische} Nachfrage ---
lebensnotwendige Güter werden auch bei Preissteigerungen kaum weniger nachgefragt, da
Ausweichmöglichkeiten fehlen. Modeartikel: $|\varepsilon| > 1$ $\Rightarrow$
\emph{elastische} Nachfrage --- auf den Artikel kann verzichtet oder ausgewichen werden,
die Nachfrager reagieren stark.

\textbf{(c)} Die Kreuzpreiselastizität misst die prozentuale Mengenänderung der Nachfrage
nach Gut~A im Verhältnis zur prozentualen Preisänderung eines anderen Gutes~B (positiv bei
Substituten, negativ bei Komplementen).
\end{loesung}

\begin{loesung}[title={Lösung Aufgabe~7 --- Preisbildung im Gleichgewicht}]
\textbf{(a)} Gleichgewichtspreis $p^* = 6$~€, Gleichgewichtsmenge $x^* = 600$ Stück
(angebotene = nachgefragte Menge).

\textbf{(b)} Bei 10~€ übersteigt das Angebot (1.000) die Nachfrage (200) um 800 Stück:
\emph{Angebotsüberschuss} (Nachfragedefizit) --- ein \emph{Käufermarkt}, der Käufer hat die
bessere Verhandlungsposition. Anbieter senken die Preise, um ihre Produktion abzusetzen;
mit sinkendem Preis steigt die Nachfrage, unrentable Anbieter drosseln die Produktion ---
der Markt bewegt sich zum Gleichgewicht.

\textbf{(c)} Bei 2~€ übersteigt die Nachfrage (1.000) das Angebot (200):
\emph{Nachfrageüberschuss} (Angebotsdefizit) --- ein \emph{Verkäufermarkt}, die Macht liegt
bei den Verkäufern. Die Anbieter können höhere Preise durchsetzen und weiten die Produktion
aus; mit steigendem Preis sinkt die Nachfrage --- Annäherung an das Gleichgewicht.

\textbf{(d)} \emph{Konsumentenrente:} Nutzengewinn der Nachfrager, die bereit gewesen
wären, mehr als den Gleichgewichtspreis zu zahlen (Differenz zwischen individueller
Zahlungsbereitschaft und $p^*$). \emph{Produzentenrente:} Mehrerlös der Anbieter, die auch
zu einem niedrigeren Preis als $p^*$ verkauft hätten.
\end{loesung}

\begin{loesung}[title={Lösung Aufgabe~8 --- Funktionen des Marktpreises}]
\emph{Signalfunktion (Informationsfunktion):} Der Preis zeigt die Knappheit eines Gutes an
(z.\,B. steigende Strompreise signalisieren Knappheit).
\emph{Ausgleichsfunktion:} Preise gleichen Angebots- und Nachfragemenge tendenziell aus
(Überschüsse werden abgebaut).
\emph{Allokationsfunktion (Lenkungsfunktion):} Preise lenken die Produktionsfaktoren dorthin,
wo marktfähige Güter entstehen (hohe Preise ziehen Investitionen an).
\emph{Selektionsfunktion (Sanktionsfunktion):} Anbieter mit zu hohen Kosten scheiden aus dem
Markt aus; Nachfrager, die den Preis nicht zahlen können oder wollen, ebenfalls.
\end{loesung}

\begin{loesung}[title={Lösung Aufgabe~9 --- Unvollständige Konkurrenz}]
\textbf{(a)} Durch \emph{Produktdifferenzierung} (Marken, Qualität, Werbung, Design) grenzt
der Anbieter sein Produkt von den Wettbewerbern ab und schafft Präferenzen bei den Kunden.
Innerhalb des so entstehenden \emph{monopolistischen Absatzbereichs} kann er den Preis in
einem Spielraum nahezu frei festlegen, ohne größere Nachfrageeinbußen. Überschreitet er den
Spielraum, wandern Kunden ab; unterschreitet er ihn, gewinnt er zwar Kunden, deckt aber
eventuell seine Kosten nicht mehr.

\textbf{(b)} \emph{Preisführerschaft:} Ein Anbieter (z.\,B. ein großer Mineralölkonzern)
erhöht vor einem Feiertag die Preise, die übrigen Tankstellenbetreiber ziehen nach.
\emph{Parallelverhalten:} Die Oligopolisten verhalten sich ohne ausdrückliche Absprache
gleichgerichtet (Preisruhe, Gleichschritt bei Preisänderungen), um Kampfsituationen zu
vermeiden; der Wettbewerb verlagert sich auf Qualität, Marke, Werbung und Service.

\textbf{(c)} Der Angebotsmonopolist ist alleiniger Anbieter und kann den Preis grundsätzlich
selbst festlegen (Preisfixierer). Grenzen: die Zahlungsbereitschaft bzw. das Ausweichen der
Nachfrager (Konsumverzicht, Substitute) sowie staatliche Regulierung und drohender
Markteintritt neuer Anbieter.
\end{loesung}

\vspace{1em}
\hrule
\vspace{0.8em}
\noindent\textsf{\small Weiterüben: Übungsband Volkswirtschaft Kap.~1 (mit Lösungsteil) und das
interaktive Quiz dieser Lektion. Nächste Lektion: Wettbewerbspolitik, Staatseingriffe \& VGR.}

\end{document}
